\documentclass[lualatex,ja=standard,a4paper,12pt]{bxjsarticle}
\usepackage{amsmath, amssymb, amsthm}
\usepackage{geometry}
\usepackage{xcolor}
\usepackage{tcolorbox}
\usepackage{booktabs}
\usepackage{array}
\usepackage{pgfplots}
\usepackage{listings}
\pgfplotsset{compat=1.18}
\usepgfplotslibrary{fillbetween}
\usetikzlibrary{patterns}

\geometry{margin=25mm}
\tcbuselibrary{theorems, skins, breakable}

\lstset{
  language=R,
  basicstyle=\ttfamily\small,
  backgroundcolor=\color{gray!10},
  frame=single,
  breaklines=true,
  keywordstyle=\color{blue},
  commentstyle=\color{green!50!black},
  stringstyle=\color{red},
}

\newtcolorbox{defbox}[1]{
  colback=blue!5, colframe=blue!60!black,
  fonttitle=\bfseries, title={#1}, arc=3pt
}
\newtcolorbox{thmbox}[1]{
  colback=green!5, colframe=green!50!black,
  fonttitle=\bfseries, title={#1}, arc=3pt,
  breakable
}
\newtcolorbox{derivbox}[1]{
  colback=yellow!6, colframe=yellow!70!black,
  fonttitle=\bfseries, title={計算補足：#1}, arc=3pt,
  breakable
}
\newtcolorbox{pointbox}{
  colback=orange!8, colframe=orange!70!black,
  fonttitle=\bfseries, title={ポイント}, arc=3pt
}
\newtcolorbox{exbox}[1]{
  colback=gray!6, colframe=gray!50,
  fonttitle=\bfseries, title={例：#1}, arc=3pt,
  breakable
}

\title{{\large Hoff 第5章} \\[6pt]
       \textbf{5.3\; Joint Inference for the Mean and Variance} \\[4pt]
       {\normalsize ゼミ資料（和訳＋計算補足）}}
\author{}
\date{}

\begin{document}
\maketitle
\thispagestyle{empty}

%===========================================
\section{問題設定：$\theta$ と $\sigma^2$ の同時推論}
%===========================================

5.2節では $\sigma^2$ を既知と仮定した．しかし現実には
\textbf{平均 $\theta$ も分散 $\sigma^2$ も両方未知}である．
5.3節ではこの状況を扱う．

目標は同時事後分布
\[
  p(\theta, \sigma^2 \mid \mathbf{y})
  \propto p(\mathbf{y} \mid \theta, \sigma^2)\, p(\theta, \sigma^2)
\]
を求め，$\theta$ および $\sigma^2$ の両方について推論することである．

%===========================================
\section{事前分布の設計}
%===========================================

二つのパラメータの同時事前分布を，条件付き確率の分解
\[
  p(\theta, \sigma^2) = p(\theta \mid \sigma^2)\, p(\sigma^2)
\]
を使って構成する．

\subsection{$\sigma^2$ の事前分布：逆ガンマ分布}

まずガンマ分布の確率密度関数を確認する．

\begin{defbox}{ガンマ分布 $\mathrm{Gamma}(a, b)$}
確率変数 $X > 0$ が shape $a > 0$，rate $b > 0$ のガンマ分布に従うとは，密度が
\[
  f_X(x) = \frac{b^a}{\Gamma(a)}\, x^{a-1}\, e^{-bx}, \quad x > 0
\]
であることを意味する（$\Gamma(a) = \int_0^\infty t^{a-1}e^{-t}dt$ はガンマ関数）．
平均 $\mathrm{E}[X] = a/b$，最頻値 $= (a-1)/b$（$a \geq 1$ のとき），分散 $= a/b^2$．
\end{defbox}

\begin{derivbox}{逆ガンマ分布の密度：変数変換による導出}
$X \sim \mathrm{Gamma}(a, b)$ のとき，$Y = 1/X$ の密度を求める．

\medskip
\noindent\textbf{変数変換の公式：}
$Y = g(X) = 1/X$（単調減少）のとき，
\[
  f_Y(y) = f_X\!\left(g^{-1}(y)\right) \cdot \left|\frac{d}{dy}g^{-1}(y)\right|.
\]
$g^{-1}(y) = 1/y$ なので $\left|\frac{d}{dy}(1/y)\right| = 1/y^2$．したがって，
\begin{align*}
  f_Y(y)
  &= f_X(1/y) \cdot \frac{1}{y^2}
   = \frac{b^a}{\Gamma(a)}\,(1/y)^{a-1}\,e^{-b/y} \cdot \frac{1}{y^2} \\[4pt]
  &= \frac{b^a}{\Gamma(a)}\, y^{-(a-1)}\, e^{-b/y} \cdot y^{-2}
   = \frac{b^a}{\Gamma(a)}\, y^{-a-1}\, e^{-b/y}.
\end{align*}

\noindent これを\textbf{逆ガンマ分布 $\mathrm{InvGamma}(a, b)$} と呼ぶ．
\end{derivbox}

\begin{defbox}{逆ガンマ分布 $\mathrm{InvGamma}(a, b)$}
$Y \sim \mathrm{InvGamma}(a, b)$ の密度：
\[
  f_Y(y) = \frac{b^a}{\Gamma(a)}\, y^{-a-1}\, e^{-b/y}, \quad y > 0.
\]
\end{defbox}

\begin{derivbox}{逆ガンマ分布のモードと平均}

\noindent\textbf{モード（密度が最大になる点）：}

$\log f_Y(y) = \mathrm{const} + (-a-1)\log y - b/y$ を $y$ で微分してゼロとおく：
\[
  \frac{d}{dy}\log f_Y(y) = \frac{-a-1}{y} + \frac{b}{y^2} = 0
  \quad\Rightarrow\quad
  (a+1)y = b
  \quad\Rightarrow\quad
  \mathrm{mode} = \frac{b}{a+1}.
\]

\noindent\textbf{平均（$a > 1$ のとき）：}

$\mathrm{E}[Y] = \mathrm{E}[1/X]$ を $X \sim \mathrm{Gamma}(a,b)$ として計算する：
\begin{align*}
  \mathrm{E}[1/X]
  &= \int_0^\infty \frac{1}{x} \cdot \frac{b^a}{\Gamma(a)}\,x^{a-1}\,e^{-bx}\,dx
   = \frac{b^a}{\Gamma(a)} \int_0^\infty x^{(a-1)-1}\,e^{-bx}\,dx \\[4pt]
  &= \frac{b^a}{\Gamma(a)} \cdot \frac{\Gamma(a-1)}{b^{a-1}}
   = \frac{b \cdot \Gamma(a-1)}{\Gamma(a)}
   = \frac{b}{a-1}.
  \quad (\Gamma(a) = (a-1)\Gamma(a-1) \text{ を使用})
\end{align*}

\noindent まとめ（$Y \sim \mathrm{InvGamma}(a,b)$）：
\[
  \mathrm{mode}(Y) = \frac{b}{a+1}, \qquad
  \mathrm{E}[Y] = \frac{b}{a-1} \quad (a > 1).
\]
\end{derivbox}

\medskip
Hoff の表記では $1/\sigma^2 \sim \mathrm{Gamma}(\nu_0/2,\; \nu_0\sigma_0^2/2)$，
すなわち $\sigma^2 \sim \mathrm{InvGamma}(\nu_0/2,\; \nu_0\sigma_0^2/2)$．

\begin{derivbox}{パラメータ $(\sigma_0^2, \nu_0)$ の意味}
$a = \nu_0/2$，$b = \nu_0\sigma_0^2/2$ を上の公式に代入する：

\begin{align*}
  \mathrm{mode}(\sigma^2)
  &= \frac{\nu_0\sigma_0^2/2}{\nu_0/2 + 1}
   = \frac{\nu_0\sigma_0^2}{\nu_0 + 2}
   \;\xrightarrow{\nu_0\to\infty}\; \sigma_0^2, \\[6pt]
  \mathrm{E}[\sigma^2]
  &= \frac{\nu_0\sigma_0^2/2}{\nu_0/2 - 1}
   = \frac{\nu_0\sigma_0^2}{\nu_0 - 2}
   \;\xrightarrow{\nu_0\to\infty}\; \sigma_0^2. \quad (\nu_0 > 2 \text{ のとき定義})
\end{align*}

\begin{itemize}
  \item $\sigma_0^2$：$\nu_0$ が大きいとき $\sigma^2$ の「中心値」に収束する．事前の「代表値」．
  \item $\nu_0$ が大きい $\Rightarrow$ モード・平均がともに $\sigma_0^2$ に近づき，分布が集中する（事前情報が強い）．
  \item $\nu_0$ が小さい $\Rightarrow$ 分布が右に広がり，$\sigma^2$ の推定が緩やかになる（事前情報が弱い）．
\end{itemize}
\end{derivbox}

\begin{derivbox}{なぜ正規モデルに対して共役か}
$\theta$ を固定して尤度を $\sigma^2$ の関数として見ると，
\[
  p(\mathbf{y} \mid \theta, \sigma^2)
  \propto (\sigma^2)^{-n/2}\, e^{-\frac{\sum(y_i-\theta)^2}{2\sigma^2}}.
\]
これは $\sigma^2$ について $\mathrm{InvGamma}$ の核 $(\sigma^2)^{-a-1}e^{-b/\sigma^2}$ と同じ形をしている
（$a = n/2 - 1$，$b = \sum(y_i-\theta)^2/2$ に対応）．

逆ガンマ事前分布 $p(\sigma^2) \propto (\sigma^2)^{-\nu_0/2-1} e^{-\nu_0\sigma_0^2/(2\sigma^2)}$ と掛け合わせると，
\[
  p(\sigma^2 \mid \mathbf{y})
  \propto (\sigma^2)^{-(\nu_0/2 + n/2)-1}\, e^{-(\nu_0\sigma_0^2 + \sum(y_i-\theta)^2)/(2\sigma^2)},
\]
これもまた逆ガンマ分布（$a$ が $\nu_0/2 + n/2$，$b$ が対応する値に更新）．
指数型の関数形が一致しているため，事前と事後が同じ族に属する＝\textbf{共役}．
\end{derivbox}

\begin{center}
\begin{tikzpicture}
\begin{axis}[
  width=11cm, height=5.5cm,
  domain=0.01:5, samples=300,
  xmin=0, xmax=5, ymin=0, ymax=1.4,
  xlabel={$\sigma^2$}, ylabel={密度},
  legend pos=north east,
  every axis/.append style={font=\small},
  title={逆ガンマ分布の形状（$\sigma_0^2=1$，$\nu_0$ を変化）},
]
% nu0=1: p(s2) prop (s2)^(-3/2)*exp(-1/(2*s2))
\addplot[blue, thick] {x^(-1.5)*exp(-1/(2*x))/(sqrt(2*pi))};
\addlegendentry{$\nu_0=1$}
% nu0=3: shape a=3/2, rate b=3/2
\addplot[red, thick] {(3/2)^(3/2)/0.886*x^(-3/2-1)*exp(-(3/2)/x)};
\addlegendentry{$\nu_0=3$}
% nu0=10: shape a=5, rate b=5
\addplot[teal, thick] {5^5/24*x^(-6)*exp(-5/x)};
\addlegendentry{$\nu_0=10$}
\end{axis}
\end{tikzpicture}
\end{center}

\subsection{$\theta$ の条件付き事前分布}

$\sigma^2$ が与えられたもとでの $\theta$ の条件付き事前分布：
\[
  \theta \mid \sigma^2 \sim \mathrm{Normal}\!\left(\mu_0,\, \frac{\sigma^2}{\kappa_0}\right).
\]

$\tau_0^2 = \sigma^2/\kappa_0$ と設定することで，$\theta$ の不確かさが $\sigma^2$ に比例するという
スケール整合性を持たせている（$\kappa_0$ は「事前の仮想観測数」）．

\bigskip
\begin{center}
\renewcommand{\arraystretch}{1.3}
\begin{tabular}{lll}
\toprule
パラメータ & 記号 & 意味 \\
\midrule
事前平均 & $\mu_0$ & $\theta$ に対する事前の予想値 \\
事前観測数（平均） & $\kappa_0$ & 小さいほど $\mu_0$ への確信が弱い \\
事前スケール & $\sigma_0^2$ & $\sigma^2$ の事前の中心値 \\
事前観測数（分散） & $\nu_0$ & 小さいほど $\sigma^2$ への確信が弱い \\
\bottomrule
\end{tabular}
\end{center}

\bigskip
モデル全体をまとめると：
\begin{align*}
  \frac{1}{\sigma^2} &\sim \mathrm{Gamma}\!\left(\frac{\nu_0}{2},\, \frac{\nu_0\sigma_0^2}{2}\right), \\[4pt]
  \theta \mid \sigma^2 &\sim \mathrm{Normal}\!\left(\mu_0,\, \frac{\sigma^2}{\kappa_0}\right), \\[4pt]
  Y_1, \ldots, Y_n \mid \theta, \sigma^2 &\overset{\mathrm{i.i.d.}}{\sim} \mathrm{Normal}(\theta,\, \sigma^2).
\end{align*}


%===========================================
\section{事後分布の導出}
%===========================================

同時事後分布を
\[
  p(\theta, \sigma^2 \mid \mathbf{y}) = p(\theta \mid \sigma^2, \mathbf{y})\, p(\sigma^2 \mid \mathbf{y})
\]
と分解して，各因子を求める．

\subsection{$\theta$ の条件付き事後分布}

$\sigma^2$ を既知として固定すると，これは 5.2節の設定に帰着する．
$\tau_0^2 = \sigma^2/\kappa_0$ に対応するので，5.2節の結果より：
\[
  \theta \mid \sigma^2, \mathbf{y}
  \;\sim\; \mathrm{Normal}\!\left(\mu_n,\, \frac{\sigma^2}{\kappa_n}\right),
\]
ただし
\[
  \kappa_n = \kappa_0 + n, \qquad
  \mu_n = \frac{\kappa_0\mu_0 + n\bar{y}}{\kappa_n}.
\]

\subsection{$\sigma^2$ の周辺事後分布：$\theta$ の積分消去}

$p(\sigma^2 \mid \mathbf{y})$ を得るには，同時事後分布から $\theta$ を積分消去する：
\[
  p(\sigma^2 \mid \mathbf{y})
  = \int_{-\infty}^{\infty} p(\theta, \sigma^2 \mid \mathbf{y})\, d\theta.
\]

\begin{derivbox}{$\theta$ の積分消去（Step by Step）}

\noindent\textbf{Step 1：同時事後分布の指数部分の整理}

\vspace{4pt}
同時事後分布の核（$\theta, \sigma^2$ 両方に依存する部分）：
\begin{align*}
  p(\theta, \sigma^2 \mid \mathbf{y})
  &\propto (\sigma^2)^{-n/2}\,
    e^{-\frac{1}{2\sigma^2}\sum(y_i-\theta)^2}
    \cdot (\sigma^2)^{-1/2}\,
    e^{-\frac{\kappa_0}{2\sigma^2}(\theta-\mu_0)^2}
    \cdot (\sigma^2)^{-\nu_0/2-1}\,
    e^{-\frac{\nu_0\sigma_0^2}{2\sigma^2}} \\[4pt]
  &= (\sigma^2)^{-(n+1)/2-\nu_0/2-1}\,
     e^{-\frac{1}{2\sigma^2}\left[\sum(y_i-\theta)^2 + \kappa_0(\theta-\mu_0)^2 + \nu_0\sigma_0^2\right]}.
\end{align*}

\noindent\textbf{Step 2：$\theta$ の二次式を平方完成}

\vspace{4pt}
$\theta$ に依存する部分だけ取り出す：
\begin{align*}
  Q(\theta)
  &= \sum(y_i-\theta)^2 + \kappa_0(\theta-\mu_0)^2 \\
  &= \underbrace{(\kappa_0 + n)}_{\kappa_n}\theta^2
     - 2\underbrace{(\kappa_0\mu_0 + n\bar{y})}_{\kappa_n\mu_n}\theta
     + \underbrace{\left(\sum y_i^2 + \kappa_0\mu_0^2\right)}_{C} \\
  &= \kappa_n(\theta - \mu_n)^2 + \underbrace{\left(C - \kappa_n\mu_n^2\right)}_{SSR_\kappa}.
\end{align*}

ここで $SSR_\kappa = \sum y_i^2 + \kappa_0\mu_0^2 - \kappa_n\mu_n^2$ とおいた（$\theta$ に依存しない）．

\noindent\textbf{Step 3：$\theta$ についてガウス積分}

\vspace{4pt}
\begin{align*}
  p(\sigma^2 \mid \mathbf{y})
  &\propto (\sigma^2)^{-(n+1)/2-\nu_0/2-1}
     e^{-\frac{SSR_\kappa + \nu_0\sigma_0^2}{2\sigma^2}}
     \int_{-\infty}^{\infty}
       e^{-\frac{\kappa_n(\theta-\mu_n)^2}{2\sigma^2}} d\theta \\[4pt]
  &= (\sigma^2)^{-(n+1)/2-\nu_0/2-1}
     e^{-\frac{SSR_\kappa + \nu_0\sigma_0^2}{2\sigma^2}}
     \cdot \sqrt{\frac{2\pi\sigma^2}{\kappa_n}} \\[4pt]
  &\propto (\sigma^2)^{-n/2-\nu_0/2-1}
     e^{-\frac{SSR_\kappa + \nu_0\sigma_0^2}{2\sigma^2}}.
\end{align*}

これは $1/\sigma^2 \sim \mathrm{Gamma}\!\left(\frac{\nu_n}{2},\, \frac{\nu_n\sigma_n^2}{2}\right)$ の核の形であり，
$\nu_n = \nu_0 + n$ かつ $\nu_n\sigma_n^2 = \nu_0\sigma_0^2 + SSR_\kappa$
を与える．
\end{derivbox}

\begin{derivbox}{$SSR_\kappa$ の整理：$\nu_n\sigma_n^2$ の公式導出}

$SSR_\kappa = \sum y_i^2 + \kappa_0\mu_0^2 - \kappa_n\mu_n^2$ を整理する．

\noindent\textbf{補題（よく使う恒等式）：} $\sum y_i^2 = (n-1)s^2 + n\bar{y}^2$

（$\sum(y_i-\bar{y})^2 = \sum y_i^2 - n\bar{y}^2$ より）

\medskip

$\kappa_n\mu_n^2 = \frac{(\kappa_0\mu_0 + n\bar{y})^2}{\kappa_n}$ に注意して，

\begin{align*}
  SSR_\kappa
  &= (n-1)s^2 + n\bar{y}^2 + \kappa_0\mu_0^2
     - \frac{(\kappa_0\mu_0 + n\bar{y})^2}{\kappa_n}.
\end{align*}

分母を $\kappa_n$ に揃えて計算：
\begin{align*}
  \kappa_n(n\bar{y}^2 + \kappa_0\mu_0^2) - (\kappa_0\mu_0 + n\bar{y})^2.
\end{align*}

$\kappa_n = \kappa_0 + n$ として展開する：
\begin{align*}
  &(\kappa_0+n)(n\bar{y}^2 + \kappa_0\mu_0^2)
   - (\kappa_0^2\mu_0^2 + 2\kappa_0 n\mu_0\bar{y} + n^2\bar{y}^2) \\
  &= \kappa_0 n\bar{y}^2 + n^2\bar{y}^2
     + \kappa_0^2\mu_0^2 + \kappa_0 n\mu_0^2
     - \kappa_0^2\mu_0^2 - 2\kappa_0 n\mu_0\bar{y} - n^2\bar{y}^2 \\
  &= \kappa_0 n(\bar{y}^2 - 2\mu_0\bar{y} + \mu_0^2)
   = \kappa_0 n(\bar{y} - \mu_0)^2.
\end{align*}

したがって：
\[
  SSR_\kappa = (n-1)s^2 + \frac{\kappa_0 n}{\kappa_n}(\bar{y} - \mu_0)^2.
\]

これを $\nu_n\sigma_n^2$ に代入：
\[
  \boxed{\nu_n\sigma_n^2
  = \nu_0\sigma_0^2 + (n-1)s^2 + \frac{\kappa_0 n}{\kappa_n}(\bar{y} - \mu_0)^2.}
\]
\end{derivbox}

\begin{thmbox}{事後分布（$\theta$ と $\sigma^2$ の同時推論）}
\[
  \theta \mid \sigma^2, \mathbf{y}
  \;\sim\; \mathrm{Normal}\!\left(\mu_n,\, \frac{\sigma^2}{\kappa_n}\right),
  \qquad
  \frac{1}{\sigma^2} \mid \mathbf{y}
  \;\sim\; \mathrm{Gamma}\!\left(\frac{\nu_n}{2},\, \frac{\nu_n\sigma_n^2}{2}\right),
\]
ただし
\begin{align*}
  \kappa_n &= \kappa_0 + n, \\
  \mu_n &= \frac{\kappa_0\mu_0 + n\bar{y}}{\kappa_n}, \\
  \nu_n &= \nu_0 + n, \\
  \nu_n\sigma_n^2 &= \nu_0\sigma_0^2 + (n-1)s^2
    + \frac{\kappa_0 n}{\kappa_n}(\bar{y} - \mu_0)^2.
\end{align*}
\end{thmbox}

\subsection{$\nu_n\sigma_n^2$ の3項の意味}

\begin{center}
\renewcommand{\arraystretch}{1.5}
\begin{tabular}{lp{8cm}}
\toprule
\textbf{項} & \textbf{意味} \\
\midrule
$\nu_0\sigma_0^2$ & 事前の「平方和」（prior sum of squares）：\\
& $\nu_0$ 個の仮想観測が$\sigma_0^2$まわりにばらついていた，という事前情報 \\[4pt]
$(n-1)s^2$ & データの平方和（sample sum of squares）：\\
& 実際の $n$ 個の観測からの $\theta$ まわりのばらつき \\[4pt]
$\dfrac{\kappa_0 n}{\kappa_n}(\bar{y}-\mu_0)^2$ & 事前平均とデータ平均のズレによる「補正項」：\\
& $\bar{y}$ と $\mu_0$ が離れているほど，$\sigma^2$ の事後推定が大きくなる \\
\bottomrule
\end{tabular}
\end{center}


%===========================================
\section{モンテカルロサンプリング}
%===========================================

$\theta$ の\textbf{周辺}事後分布 $p(\theta \mid \mathbf{y})$ は，
$p(\theta \mid \sigma^2, \mathbf{y})$ と $p(\sigma^2 \mid \mathbf{y})$ の積を $\sigma^2$ について積分したものだが，
解析的に計算するのは煩雑である．そこで\textbf{モンテカルロ法}を用いる．

\begin{thmbox}{逐次サンプリング（Ancestral Sampling）}
次の手順を $S$ 回繰り返す（$s = 1, \ldots, S$）：
\begin{enumerate}
  \item $\sigma^{2(s)} \sim \mathrm{InvGamma}\!\left(\dfrac{\nu_n}{2},\, \dfrac{\nu_n\sigma_n^2}{2}\right)$
        \hfill（$1/\sigma^{2(s)} \sim \mathrm{Gamma}(\nu_n/2,\, \nu_n\sigma_n^2/2)$ とサンプル）
  \item $\theta^{(s)} \sim \mathrm{Normal}\!\left(\mu_n,\, \dfrac{\sigma^{2(s)}}{\kappa_n}\right)$
\end{enumerate}

得られた $\{(\theta^{(1)}, \sigma^{2(1)}), \ldots, (\theta^{(S)}, \sigma^{2(S)})\}$ は
同時事後分布 $p(\theta, \sigma^2 \mid \mathbf{y})$ からの独立なサンプルである．
\end{thmbox}

\begin{derivbox}{なぜこれが周辺分布のサンプルになるのか}

$(\theta^{(s)}, \sigma^{2(s)})$ の同時分布を確認する：

\medskip
\noindent Step 1 により $\sigma^{2(s)} \sim p(\sigma^2 \mid \mathbf{y})$.

\noindent Step 2 により $\theta^{(s)} \mid \sigma^{2(s)} \sim p(\theta \mid \sigma^{2(s)}, \mathbf{y})$.

\medskip
これは「$(\theta, \sigma^2)$ を $p(\theta, \sigma^2 \mid \mathbf{y})$ からサンプルする」ことと全く同じである：
\[
  p(\theta^{(s)}, \sigma^{2(s)})
  = p(\theta^{(s)} \mid \sigma^{2(s)}, \mathbf{y})\, p(\sigma^{2(s)} \mid \mathbf{y})
  = p(\theta^{(s)}, \sigma^{2(s)} \mid \mathbf{y}).
\]

そのため，$\theta^{(s)}$ の（$\sigma^{2(s)}$ を問わない）\textbf{周辺分布}は：
\[
  p(\theta^{(s)}) = \int p(\theta^{(s)}, \sigma^{2(s)} \mid \mathbf{y})\, d\sigma^{2(s)}
  = p(\theta^{(s)} \mid \mathbf{y}).
\]

直感：各 $\theta^{(s)}$ は「異なる $\sigma^{2(s)}$ のもとで生成されている」．
集めると，その生成に使った $\sigma^{2(s)}$ の分布（$= p(\sigma^2|\mathbf{y})$）にわたって平均しているので，
結果的に $\sigma^2$ を積分消去したことになる．
\end{derivbox}

\noindent R コードでは 3 行で実装できる：

\begin{lstlisting}
# 事後パラメータを設定後...
phi.post  <- rgamma(S, shape = nun/2, rate = nun*s2n/2)   # 精度のサンプル
s2.post   <- 1/phi.post                                   # 分散に変換
theta.post <- rnorm(S, mean = mun, sd = sqrt(s2.post/kn)) # 平均のサンプル
\end{lstlisting}

\begin{lstlisting}
# 95% 信用区間（周辺）
quantile(theta.post, c(0.025, 0.975))
quantile(s2.post, c(0.025, 0.975))
\end{lstlisting}


%===========================================
\section{無情報事前分布との接続：$t$ 分布の出現}
%===========================================

\subsection{無情報極限}

$\kappa_0, \nu_0 \to 0$（事前情報をゼロに近づける）とすると：
\begin{align*}
  \mu_n &\to \bar{y}, \\
  \kappa_n &\to n, \\
  \nu_n &\to n, \\
  \nu_n\sigma_n^2 &\to (n-1)s^2,
  \quad \text{すなわち} \quad
  \sigma_n^2 \to \frac{(n-1)s^2}{n}.
\end{align*}

\subsection{$\theta$ の周辺事後分布は $t$ 分布}

\begin{derivbox}{周辺事後分布の計算（ガンマ積分）}
$\theta|\mathbf{y}$ の周辺事後分布は，$\sigma^2$ を積分消去することで得られる：
\[
  p(\theta \mid \mathbf{y})
  = \int_0^\infty p(\theta \mid \sigma^2, \mathbf{y})\, p(\sigma^2 \mid \mathbf{y})\, d\sigma^2.
\]

$\phi = 1/\sigma^2$ とおくと（$\phi \mid \mathbf{y} \sim \mathrm{Gamma}(\nu_n/2, \nu_n\sigma_n^2/2)$）：
\begin{align*}
  p(\theta \mid \mathbf{y})
  &\propto \int_0^\infty \phi^{1/2}\, e^{-\frac{\kappa_n\phi}{2}(\theta-\mu_n)^2}
    \cdot \phi^{\nu_n/2-1}\, e^{-\frac{\nu_n\sigma_n^2}{2}\phi}\, d\phi \\
  &= \int_0^\infty \phi^{(\nu_n+1)/2-1}\,
     e^{-\frac{\phi}{2}\left[\kappa_n(\theta-\mu_n)^2 + \nu_n\sigma_n^2\right]}\, d\phi.
\end{align*}

これはガンマ関数の積分公式
$\displaystyle\int_0^\infty \phi^{a-1} e^{-b\phi}\, d\phi = \Gamma(a)/b^a$
を $a = (\nu_n+1)/2$，$b = \left[\kappa_n(\theta-\mu_n)^2 + \nu_n\sigma_n^2\right]/2$ で適用する：
\begin{align*}
  p(\theta \mid \mathbf{y})
  &\propto \left(\kappa_n(\theta-\mu_n)^2 + \nu_n\sigma_n^2\right)^{-(\nu_n+1)/2} \\
  &= \left(\nu_n\sigma_n^2\right)^{-(\nu_n+1)/2}
     \left(1 + \frac{\kappa_n(\theta-\mu_n)^2}{\nu_n\sigma_n^2}\right)^{-(\nu_n+1)/2} \\
  &\propto \left(1 + \frac{1}{\nu_n}
     \left(\frac{\theta-\mu_n}{\sqrt{\sigma_n^2/\kappa_n}}\right)^2
    \right)^{-(\nu_n+1)/2}.
\end{align*}

これは自由度 $\nu_n$，位置 $\mu_n$，スケール $\sqrt{\sigma_n^2/\kappa_n}$ の $t$ 分布の核である．
\end{derivbox}

\begin{thmbox}{$\theta$ の周辺事後分布（$t$ 分布）}
\[
  \frac{\theta - \mu_n}{\sqrt{\sigma_n^2/\kappa_n}} \;\Bigg|\; \mathbf{y}
  \;\sim\; t_{\nu_n}.
\]
無情報事前分布の極限（$\kappa_0, \nu_0 \to 0$）では
$\mu_n \to \bar{y}$，$\nu_n \to n$，$\sigma_n^2/\kappa_n \to s^2/n$ となり，
\[
  \frac{\theta - \bar{y}}{s/\sqrt{n}} \;\Bigg|\; \mathbf{y}
  \;\sim\; t_{n-1}.
\]
\end{thmbox}

これは頻度論の $t$ 検定と\textbf{数式として同じ形}をしているが，解釈が異なる：

\begin{center}
\renewcommand{\arraystretch}{1.5}
\begin{tabular}{lp{5cm}p{5.5cm}}
\toprule
 & \textbf{式} & \textbf{解釈} \\
\midrule
頻度論 &
  $\dfrac{\bar{Y} - \theta}{s/\sqrt{n}} \sim t_{n-1}$ &
  $\theta$ は固定定数，$\bar{Y}$ がデータ間でランダムに変動する \\[8pt]
ベイズ（無情報） &
  $\dfrac{\theta - \bar{y}}{s/\sqrt{n}} \sim t_{n-1}$ &
  $\bar{y}$ は固定観測値，$\theta$ が事後分布に従いランダムに変動する \\
\bottomrule
\end{tabular}
\end{center}

\begin{pointbox}
95\% 信頼区間（頻度論）と 95\% 信用区間（ベイズ，無情報事前）は
\textbf{数値的に一致}するが意味が異なる：
\begin{itemize}
  \item 信頼区間：「データを繰り返し取り直すと 95\% の割合で $\theta$ を含む区間」
  \item 信用区間：「観測されたデータのもとで $\theta$ がこの区間にいる確率が 95\%」（直接的な確率文）
\end{itemize}
\end{pointbox}


%===========================================
\section{例：ユスリカの翅長データ（$\sigma^2$ 未知）}
%===========================================

\begin{exbox}{5.2節の再訪（$\sigma^2$ を未知として同時推論）}
5.2節と同じデータ：$n=9$，$\bar{y} = 1.804$，$s^2 = 0.017$

\medskip
事前分布のパラメータを次のように設定する：
\[
  \mu_0 = 1.9, \quad \kappa_0 = 1, \quad \sigma_0^2 = 0.01, \quad \nu_0 = 1.
\]
\begin{itemize}
  \item $\kappa_0 = 1$：事前情報は「1 個の仮想観測」程度（弱い）
  \item $\nu_0 = 1$，$\sigma_0^2 = 0.01$：$\sigma^2$ は $0.01$ 程度と予想するが不確か
\end{itemize}
\end{exbox}

\begin{derivbox}{事後パラメータの数値計算}
\begin{align*}
  \kappa_n &= 1 + 9 = 10, \\[4pt]
  \mu_n &= \frac{1 \times 1.9 + 9 \times 1.804}{10} = \frac{18.136}{10} = 1.814, \\[4pt]
  \nu_n &= 1 + 9 = 10, \\[4pt]
  \nu_n\sigma_n^2
  &= \underbrace{1 \times 0.01}_{\nu_0\sigma_0^2}
   + \underbrace{8 \times 0.017}_{(n-1)s^2}
   + \underbrace{\frac{1 \times 9}{10}(1.804 - 1.9)^2}_{\text{補正項}} \\
  &= 0.010 + 0.136 + 0.9 \times 0.009216 \\
  &= 0.010 + 0.136 + 0.0083 = 0.154, \\[4pt]
  \sigma_n^2 &= 0.154/10 = 0.0154.
\end{align*}

周辺事後分布：
$\dfrac{\theta - 1.814}{\sqrt{0.0154/10}} \mid \mathbf{y} \sim t_{10}$，
すなわち $\dfrac{\theta - 1.814}{0.0392} \mid \mathbf{y} \sim t_{10}$.

95\% 信用区間（$t_{10}$ の 2.5\% 点 $\approx \pm 2.228$）：
\[
  1.814 \pm 2.228 \times 0.0392 \approx (1.727,\; 1.901).
\]

\noindent\textbf{5.2節（$\sigma^2 = s^2$ 既知と仮定）の 95\% 区間 $(1.72,\; 1.89)$ との比較：}
\begin{itemize}
  \item $\sigma^2$ 未知の方が区間が少し広い（$\sigma^2$ の不確かさ分）
  \item ただし $n=9$ と小さいわりには差は小さい（$\kappa_0, \nu_0$ が小さいため）
\end{itemize}
\end{derivbox}

\begin{center}
\begin{tikzpicture}
\begin{axis}[
  width=12cm, height=6cm,
  domain=1.5:2.1, samples=400,
  xmin=1.5, xmax=2.1, ymin=0, ymax=12,
  xlabel={$\theta$（翅長の母平均，mm）},
  ylabel={密度},
  legend pos=north east,
  every axis/.append style={font=\small},
]
% t_{10}(1.814, 0.0154/10): scale = sqrt(0.0154/10) = 0.0392
% t density: C*(1 + ((x-mu)/s)^2 / nu)^(-(nu+1)/2)
% C = Gamma(5.5)/(sqrt(10*pi)*Gamma(5)) = 6.5625/(24*sqrt(10*pi)) approx 0.4869
% Using nu=10, mu=1.814, s=0.0392
\addplot[red, thick]
  {0.4869/0.0392 * (1 + ((x-1.814)/0.0392)^2/10)^(-5.5)};
\addlegendentry{$t_{10}(1.814, 0.0392^2)$（$\sigma^2$ 未知の事後）}
% Normal(1.805, 0.002): sigma = sqrt(0.002) = 0.0447
\addplot[blue, thick, dashed]
  {1/(0.0447*sqrt(2*pi))*exp(-((x-1.805)^2)/(2*0.002))};
\addlegendentry{$\mathrm{Normal}(1.805, 0.002)$（5.2節：$\sigma^2=0.017$ 既知）}
\end{axis}
\end{tikzpicture}
\end{center}

\noindent $\sigma^2$ 未知の事後分布（$t$ 分布）は，既知の場合（正規分布）よりも裾が重く，
$\sigma^2$ の不確かさが $\theta$ の推論に伝播していることが見てとれる．


%===========================================
\section{まとめ}
%===========================================

\begin{center}
\renewcommand{\arraystretch}{1.4}
\begin{tabular}{lp{9.5cm}}
\toprule
\textbf{項目} & \textbf{内容} \\
\midrule
事前分布 &
  $p(\theta, \sigma^2) = p(\theta|\sigma^2) \cdot p(\sigma^2)$ に分解．\\
  & $1/\sigma^2 \sim \text{Gamma}(\nu_0/2, \nu_0\sigma_0^2/2)$，
    $\theta|\sigma^2 \sim N(\mu_0, \sigma^2/\kappa_0)$ \\
共役性 &
  事後分布も同じ形（正規-逆ガンマ族）に留まる \\
$\sigma_n^2$ の公式 &
  3項の和：事前平方和 $+$ データ平方和 $+$ ズレのペナルティ \\
MC サンプリング &
  $\sigma^2 \to \theta$ の順にサンプルするだけで周辺分布が得られる \\
無情報極限 &
  $\kappa_0, \nu_0 \to 0$ で頻度論の $t_{n-1}$ 検定と数値的に一致 \\
$\sigma^2$ 未知の効果 &
  $\sigma^2$ の不確かさが裾の重さ（自由度 $\nu_n$）として伝播する \\
\bottomrule
\end{tabular}
\end{center}

\end{document}
